\clearpage
\section{Background}

This chapter establishes the technical and conceptual grounding the work rests on. Section~2.1 covers the foundations: component-based software engineering in embedded systems, the dependency injection and inversion of control patterns that the \textit{diff} framework uses to compose systems, the \textit{diff} framework itself, and the module distribution path from which it draws components. Section~2.2 reviews related work in visual composition for software systems and in cross-platform desktop application frameworks.

\subsection{Foundations}

\subsubsection{Component-Based Software Engineering in Embedded Systems}

Component-Based Software Engineering (CBSE) is the practice of building software from independent units called components, each with explicit interfaces that declare what it provides and what it requires from other components~\cite{szyperski2002}.

In an embedded system, a component is a unit like a network link driver, a timer, or a message source. Embedded products usually ship as a family of related devices that share most of their software but differ in hardware~\cite{clements2001spl}. Writing a separate codebase per device does not scale. Every shared bug fix has to be applied to every codebase, and the maintenance cost grows with the number of devices in the family. CBSE supports this setting because a component, once written, can appear in any device that needs it, and each device's software is composed from a different subset of the available components~\cite{crnkovic2002}.

CBSE on its own leaves two questions open. The first, how a component declares the other components it requires, is answered by dependency injection (\S\ref{sec:di-ioc}). The second, how the assembled system resolves those declarations into concrete instances, is answered by the \textit{diff} runtime (\S\ref{sec:diff-framework}).

\subsubsection{Dependency Injection and Inversion of Control}
\label{sec:di-ioc}

Inversion of Control (IoC) is the design principle in which a framework, rather than the component itself, decides when and how the component is instantiated and connected to its dependencies~\cite{fowler2004di}. Dependency Injection (DI) implements IoC by separating the component from the construction of its dependencies: a component declares what it requires, and an external assembler supplies it, typically through constructor parameters or setter methods.

DI changes the direction of control. Without DI, a component creates its dependencies internally:
\begin{lstlisting}[language=C++, aboveskip=10pt]
class Component {
    Database db{DatabaseFactory::create()};  // self-created
};
\end{lstlisting}

With DI, the dependencies are passed in by an external assembler:
\begin{lstlisting}[language=C++, aboveskip=10pt]
class Component {
  public:
      Component(Database& db);  // injected via constructor
  private:
      Database& db_;
};
\end{lstlisting}

The same \texttt{Component} class can take a real \texttt{Database} in production, a mock one in tests, and a different implementation when wired into a separate product variant.

\paragraph{Static vs.\ Runtime DI}
DI has two main forms, distinguished by when the wiring is resolved (Figure~\ref{fig:static-vs-runtime-di}).

\begin{figure}[ht]
    \centering
    \resizebox{\textwidth}{!}{
        \begin{tikzpicture}[
    node distance=17mm and 10mm,
    box/.style={draw, rounded corners=2pt, minimum height=10mm, text width=24mm, align=center, font=\small},
    artifact/.style={box, fill=gray!10},
    process/.style={box, fill=blue!10},
    wiring/.style={box, fill=blue!10, very thick, draw=orange!80!black},
    output/.style={box, fill=green!10},
    arr/.style={->, thick},
    arrwiring/.style={->, very thick, draw=orange!80!black},
    panel/.style={font=\small\bfseries}
]
% --- Panel A: Static DI ---
\node[panel] (titleA) {Static DI};
\node[artifact, below=4mm of titleA.south west, anchor=north west] (srcA) {Source code\\+ headers};
\node[wiring, right=of srcA] (compA) {Compiler\\(links graph)};
\node[output, right=of compA] (binA) {Wired binary};
\node[output, right=of binA] (sysA) {Live system};
\draw[arr] (srcA) -- (compA);
\draw[arr] (compA) -- (binA);
\draw[arr] (binA) -- (sysA);
\node[font=\scriptsize\itshape, color=orange!80!black, below=0.5mm of compA] {wiring resolved};

% --- Panel B: Runtime DI ---
\node[panel, below=16mm of srcA.south west, anchor=north west] (titleB) {Runtime DI};
\node[artifact, below=4mm of titleB.south west, anchor=north west] (srcB) {Source code};
\node[process, right=of srcB] (compB) {Compiler};
\node[output, right=of compB] (binB) {Generic binary};
\node[output, right=of binB] (sysB) {Live system};
\node[artifact, below=of binB] (jsonB) {JSON\\description};
\draw[arr] (srcB) -- (compB);
\draw[arr] (compB) -- (binB);
\draw[arrwiring] (binB) -- node[above, font=\scriptsize\itshape, color=orange!80!black]{loader} (sysB);
\draw[arrwiring] (jsonB) -- node[midway, sloped, above, font=\scriptsize\itshape, color=orange!80!black]{wiring resolved} (sysB);
\end{tikzpicture}

    }
    \caption{Static DI vs.\ runtime DI.}
    \label{fig:static-vs-runtime-di}
\end{figure}

\textbf{Static DI} resolves wiring at compile time: a header or template specialisation names the concrete type for each interface, the compiler links the graph, and the resulting binary has no wiring step at startup. The compiler rejects unfilled dependencies, so the system is verifiable. The binary contains exactly one topology, so every variant requires a rebuild.

\textbf{Runtime DI} resolves wiring at process startup from a data description (XML, JSON, YAML). The same binary takes different descriptions and produces different systems, so one build serves multiple variants. The compiler no longer checks the wiring, so a malformed description fails only when the runtime loader executes it.

\paragraph{DI in C++14}
relies on template metaprogramming and static initialisation. C++ has no reflection-based DI container like Java or C\#, so frameworks build their own using the Curiously Recurring Template Pattern~\cite{coplien1995}. C++14~\cite{cpp14standard} adds variadic templates and stronger \texttt{constexpr} support that simplify this construction. Niespodziany~\cite{niespodziany2025} details how the \textit{diff} framework (\S\ref{sec:diff-framework}) applies these techniques.

\FloatBarrier

\subsubsection{The \textit{diff} Framework}
\label{sec:diff-framework}

The \textit{diff} framework~\cite{niespodziany2025} is a runtime DI implementation for C++14 embedded software. Figure~\ref{fig:diff-workflow} sketches its end-to-end workflow: component classes are compiled into the binary and register themselves before \texttt{main()} runs, then a JSON topology file drives the construction of a live system. The examples below come from \textit{diff\_broker}, the reference application that ships with the framework.

\begin{figure}[ht]
    \centering
    \resizebox{\textwidth}{!}{
        \begin{tikzpicture}[
    node distance=6mm and 8mm,
    box/.style={draw, rounded corners=2pt, minimum height=9mm, text width=34mm, align=center, font=\small, inner sep=2pt},
    source/.style={box, fill=gray!10, text width=24mm},
    registry/.style={box, fill=orange!20, draw=orange!80!black, thick},
    process/.style={box, fill=blue!10},
    output/.style={box, fill=green!10, text width=28mm},
    arr/.style={->, thick},
    dasharr/.style={->, thick, dashed, draw=orange!80!black},
    phase/.style={font=\small\bfseries\itshape, anchor=west},
    every node/.append style={execute at begin node={\hyphenpenalty=10000\exhyphenpenalty=10000}}
]
% --- Phase 1: static-init (before main) ---
\node[phase] (phaseA) at (0,0) {Compile + static-init};

\node[source, below=2mm of phaseA.south west, anchor=north west] (cls) {Component\\classes};
\node[registry, right=22mm of cls] (factreg) {FactoryRegistry};
\draw[arr] (cls) -- node[above, font=\scriptsize\itshape]{registered} (factreg);

% --- Phase 2: runtime (in main) ---
\node[phase, below=16mm of phaseA.south west, anchor=north west] (phaseB) {Runtime};

\node[source, below=2mm of phaseB.south west, anchor=north west] (topo) {topology.json};
\node[process, right=of topo] (loader) {TopologyLoader};
\node[process, right=of loader] (build) {\texttt{diff::Build}};
\node[registry, right=of build] (depreg) {DependencyRegistry};
\node[output, below=8mm of depreg] (app) {Application\\\texttt{\scriptsize build.get<I>()}};

\draw[arr] (topo) -- (loader);
\draw[arr] (loader) -- (build);
\draw[arr] (build) -- (depreg);
\draw[arr] (depreg) -- (app);

% --- Cross-phase: FactoryRegistry feeds Build ---
\draw[dasharr] (factreg.south) -- ++(0,-8mm) -| (build.north) node[pos=0.25, above, font=\scriptsize\itshape, color=orange!80!black] {lookup by type};
\end{tikzpicture}

    }
    \caption{The \textit{diff} workflow. Component classes register their factories at static-initialisation time. At runtime, the framework reads the topology JSON, looks up the factories by type name, and instantiates the components into a live application accessible by interface.}
    \label{fig:diff-workflow}
\end{figure}

\paragraph{Components and interface components}
Niespodziany~\cite{niespodziany2025} distinguishes two kinds of components. A regular \textit{component} implements algorithms and data structures. An \textit{interface component} is header-only: a pure-virtual C++ class that declares an interface and lives in its own translation unit so no implementation can leak detail across the contract. A component class inherits from \texttt{diff::Component<T, ...>}, where \texttt{T} is the class itself and the remaining template parameters declare which interfaces the component provides. Listing~\ref{lst:diff-component} shows \texttt{LinkGsm}: it implements \texttt{ILink} and exposes itself under that interface, so any later component that needs an \texttt{ILink} can be wired to it.

\begin{addmargin}[0mm]{0mm}
\begin{lstlisting}[language=C++, caption={A \textit{diff} component that implements the \texttt{ILink} interface.}, label=lst:diff-component, aboveskip=10pt]
// ILink.h -- interface component (header-only)
class ILink {
public:
    virtual ~ILink() = default;
    virtual bool send(const std::string& message) = 0;
};

// LinkGsm.h
class LinkGsm : public diff::Component<LinkGsm, diff::as<ILink>> {
public:
    LinkGsm();
    bool send(const std::string& message) override;
};
\end{lstlisting}
\end{addmargin}

\paragraph{The topology file}
The system's wiring is described by a JSON \textit{topology} file. Each entry carries the component's type name, a unique instance id, the ordered ids of the dependencies that should be injected into the constructor, and a typed configuration map (Listing~\ref{lst:topology-entry}). The framework instantiates the entries in order, so the topology must be in dependency order: every component appears after all the components it depends on.

\begin{addmargin}[0mm]{0mm}
\begin{lstlisting}[caption={Topology entry consumed by the \textit{diff} loader at startup.}, label=lst:topology-entry, aboveskip=10pt]
{
  "type": "MessageSource",
  "id": "msgSrc0",
  "version": "1.0.0",
  "source": "https://artifactory.example.com/artifactory/diff",
  "dependencies": ["linkEth0", "linkGsm0"],
  "config": {
    "count": 3,
    "content": "Diff says hello."
  }
}
\end{lstlisting}
\end{addmargin}

\paragraph{Topology as a data product}
The same compiled binary, given two different topology files, produces two different running systems with different components and different wiring. This is what Niespodziany~\cite{niespodziany2025} calls \textit{data-driven} composition. A new variant is a JSON edit, not a code change. The trade-off is that the JSON file becomes the primary artifact of the system, and any error in it shows up only when the loader runs. That trade-off is what motivates the visual editor presented in this thesis.

\subsubsection{Module Distribution: Conan and Artifactory}

C++ has no standard package manager, so teams pick one of several third-party tools to manage binary dependencies. \textbf{Conan}~\cite{conan} is the one used by the \textit{diff} ecosystem: a package manager that builds and consumes C++ libraries through \textit{recipes}, build scripts that describe how to compile, package, and install a library. A Conan package is identified by \texttt{name/version@user/channel} and may carry pre-built binaries for multiple compilers, architectures, and build types.

\textbf{Artifactory}~\cite{artifactory} is a binary repository manager from JFrog that, among other formats, supports the Conan protocol. Teams publish their internal C++ libraries as Conan packages to an Artifactory repository, and consumer projects fetch them through standard Conan tooling. In a \textit{diff}-based environment, each component is a Conan package, and each package carries a catalog entry that describes the component to potential consumers: the interfaces it implements, the dependencies it requires, and the configuration parameters it accepts. Listing~\ref{lst:catalog-entry} shows the catalog entry for a \texttt{MessageSource} component.

\begin{addmargin}[0mm]{0mm}
\begin{lstlisting}[caption={Catalog entry for the \texttt{MessageSource} component, published alongside its Conan package.}, label=lst:catalog-entry, aboveskip=10pt]
{
  "type": "MessageSource",
  "version": "1.0.0",
  "implements": ["IProcessable"],
  "requires": [
    { "slot": "link",    "type": "ILink" },
    { "slot": "backup",  "type": "ILink", "isArray": true }
  ],
  "config": {
    "count":   { "type": "uint32", "min": 1, "max": 1000 },
    "content": { "type": "string", "default": "msg" }
  }
}
\end{lstlisting}
\end{addmargin}

The catalog entry is a \textit{schema} rather than an instance: \texttt{implements} lists the interfaces the component provides, \texttt{requires} describes the slots that must be filled by other components (each with a name, an interface type, and an \texttt{isArray} flag for multi-binding slots), and \texttt{config} declares each parameter's type, allowed range, and default value. A topology entry (\S\ref{sec:diff-framework}) is a concrete \textit{instance} of a catalog entry: it picks an \texttt{id}, fills each \texttt{requires} slot with the \texttt{id} of another instance, and assigns concrete values inside \texttt{config}.

For \textit{diff-forge}, the catalog of available modules is the set of catalog entries hosted on Artifactory. At startup, the tool queries the repository's Conan API, downloads each component's catalog entry, and presents the entries in the side panel. The ingestion pipeline is described in \S4.3.

\subsection{Related Work}

\subsubsection{Visual Composition for Software Systems}

Visual composition tools let developers build software by arranging and connecting nodes on a canvas, replacing or supplementing textual code. The approach has a long history in domain-specific environments. LabVIEW~\cite{labview} uses a dataflow paradigm for instrumentation and measurement software. Simulink~\cite{simulink} applies it to control-system modelling. Node-RED~\cite{nodered} offers a flow-based editor for IoT devices and web services. Unreal Engine's Blueprint system~\cite{blueprint} exposes game-engine APIs as a visual scripting language.

In all of these tools, the canvas is a directed graph. Nodes represent units of computation or component instances. Edges represent dataflow or dependency. The two graph-theory concepts relevant to \textit{diff-forge} are the directed acyclic graph and the algorithm for ordering its nodes.

\paragraph{Directed Acyclic Graph (DAG)}
is a finite directed graph that contains no directed cycles~\cite{cormen2009}. Each edge has a direction, and no sequence of edges returns to its starting node. DAGs are the natural model for dependency relationships in software. Each component depends on zero or more others, and a cycle would mean a component transitively depends on itself. A topology with a cycle cannot be instantiated, because no component in the cycle has all of its dependencies available.

The \textit{diff} framework requires the topology to be a DAG. The loader refuses to instantiate a topology that contains a cycle. \textit{diff-forge} enforces the DAG property at export time. When the user attempts to export a topology that contains a cycle, the validation engine rejects it and points at the cycle's edges.

\paragraph{Topological Sorting}
produces a linear ordering of a DAG's nodes such that for every directed edge $(u, v)$, $u$ comes before $v$~\cite{cormen2009}. For a dependency graph, a topological order is a valid construction sequence. Every component's dependencies appear before the component itself.

Kahn's algorithm~\cite{kahn1962} computes a topological order in linear time (Figure~\ref{fig:kahn-algorithm}). It begins with the set of nodes that have no incoming edges. At each step, it removes one such node from the set, appends it to the output, and reduces the in-degree of each neighbour. Any neighbour that drops to zero in-degree joins the set. The algorithm terminates when the set is empty. If all nodes have been emitted, the output is a valid topological order. If some nodes remain unemitted, the graph contains a cycle.

\begin{figure}[ht]
    \centering
    \resizebox{\textwidth}{!}{
        \begin{tikzpicture}[
    node/.style={circle, draw, thick, minimum size=0.8cm, fill=white},
    arrow/.style={->, thick},
    stepbox/.style={rectangle, draw, dashed, inner sep=0.3cm, fill=gray!5}
]

% Step 1
\begin{scope}[local bounding box=s1]
    \node[node, fill=green!10] (1a) {1};
    \node[node, fill=green!10, right=1cm of 1a] (2a) {2};
    \node[node, below=0.8cm of 1a] (3a) {3};
    \node[node, right=1cm of 3a] (4a) {4};
    \draw[arrow] (1a) -- (3a);
    \draw[arrow] (2a) -- (3a);
    \draw[arrow] (3a) -- (4a);
    \node[below=0.1cm of 3a, xshift=0.5cm] {\small \textbf{Step 1:} Find nodes with in-degree 0 (1, 2)};
\end{scope}

% Step 2
\begin{scope}[shift={(6,0)}, local bounding box=s2]
    \node[node, opacity=0.3] (1b) {1};
    \node[node, opacity=0.3, right=1cm of 1b] (2b) {2};
    \node[node, fill=green!10, below=0.8cm of 1b] (3b) {3};
    \node[node, right=1cm of 3b] (4b) {4};
    \draw[arrow, opacity=0.3] (1b) -- (3b);
    \draw[arrow, opacity=0.3] (2b) -- (3b);
    \draw[arrow] (3b) -- (4b);
    \node[below=0.1cm of 3b, xshift=0.5cm] {\small \textbf{Step 2:} Remove 1, 2. Node 3 in-degree $\to$ 0};
\end{scope}

% Step 3
\begin{scope}[shift={(12,0)}, local bounding box=s3]
    \node[node, opacity=0.3] (1c) {1};
    \node[node, opacity=0.3, right=1cm of 1c] (2c) {2};
    \node[node, opacity=0.3, below=0.8cm of 1c] (3c) {3};
    \node[node, fill=green!10, right=1cm of 3c] (4c) {4};
    \draw[arrow, opacity=0.3] (1c) -- (3c);
    \draw[arrow, opacity=0.3] (2c) -- (3c);
    \draw[arrow, opacity=0.3] (3c) -- (4c);
    \node[below=0.1cm of 3c, xshift=0.5cm] {\small \textbf{Step 3:} Remove 3. Node 4 in-degree $\to$ 0};
\end{scope}

% Final Result
\node[below=1.5cm of s2, align=center, fill=yellow!10, draw, thick] {
    \textbf{Final Topological Order:}\\
    $[1, 2, 3, 4]$
};

\end{tikzpicture}

    }
    \caption{Step-by-step topological sorting using Kahn's algorithm.}
    \label{fig:kahn-algorithm}
\end{figure}

\textit{diff-forge} uses Kahn's algorithm during export. The topology JSON file requires entries to appear in load-compatible order (leaves before roots). Kahn's algorithm produces this order from the canvas graph (\S4.5).

\subsubsection{Cross-Platform Desktop Application Frameworks}

A cross-platform desktop application framework lets a single codebase target Windows, macOS, and Linux without per-platform rewrites. The candidates differ along three dimensions. The first is the rendering technology, either a web stack or native widgets. The second is the language used for application logic. The third is the deployment profile, such as binary size and memory footprint. Table~\ref{tab:cross-platform} summarises the most active frameworks at the time of writing.

\begin{table}[ht]
\centering
\caption{Cross-platform desktop application frameworks compared. Figures are typical for a minimal application bundled for production distribution.}
\label{tab:cross-platform}
\footnotesize
\begin{tabular}{lllll}
\toprule
\textbf{Framework} & \textbf{Renderer} & \textbf{App Language} & \textbf{Binary Size} & \textbf{Memory Footprint} \\
\midrule
Electron       & Chromium + Node.js & TypeScript / JavaScript & 100--200 MB & 200--500 MB \\
Tauri          & Native WebView     & Rust + TypeScript       & 10--30 MB   & 50--150 MB  \\
Qt (Widgets)   & Native widgets     & C++                     & 40--100 MB  & 50--200 MB  \\
.NET MAUI      & Native + WebView   & C\# / XAML              & 50--100 MB  & 100--300 MB \\
Flutter Desktop & Skia              & Dart                    & 30--60 MB   & 80--200 MB  \\
\bottomrule
\end{tabular}
\end{table}

\textbf{Electron}~\cite{electronjs} bundles a Chromium browser and a Node.js runtime. The application UI is written as a standard web page, and the Node.js backend provides access to OS-level APIs. The two-process architecture (main process plus renderer process) and the IPC layer between them are described in \S5.2. The trade-off is binary size and memory footprint, both larger than native alternatives.

\textbf{Tauri}~\cite{tauri} uses the OS-provided WebView (WebKit on macOS, WebView2 on Windows, WebKitGTK on Linux) for the UI and Rust for the backend. This achieves an order-of-magnitude reduction in binary size and memory at the cost of a less uniform rendering surface across platforms.

\textbf{Qt}~\cite{qt} is the long-standing C++ framework for native desktop applications. The Widgets module produces compact, native-looking applications. The QML/Quick module offers a declarative UI with hardware acceleration. The trade-off is that the application logic is written in C++, which limits the available ecosystem of UI libraries compared to the web stack.

\textbf{.NET MAUI}~\cite{maui} is Microsoft's cross-platform framework. It combines native widgets and embedded WebViews under a C\#/XAML programming model. Its primary target is mobile (Android and iOS) with desktop support as a secondary use case.

\textbf{Flutter Desktop}~\cite{flutter} uses Google's Skia renderer to draw all UI elements directly. The result is a consistent look across platforms at the cost of not matching native conventions. Its primary language is Dart.

For \textit{diff-forge}, Electron with React and TypeScript was selected (\S4.1). Familiarity with the web stack and the maturity of the React Flow canvas library~\cite{react_flow} outweighed the binary-size cost.
